\begin{frame}[t]{2 formas de executar um algoritmo}
\begin{columns}[T]
\begin{column}{0.48\textwidth}
\begin{block}{Programa de cálculo}
Recebe duas leituras, calcula a média, exibe o resultado e termina.
\end{block}
\end{column}
\begin{column}{0.48\textwidth}
\begin{block}{Programa de controle}
Lê entradas, aplica a regra, atualiza saídas e repete enquanto estiver em funcionamento.
\end{block}
\end{column}
\end{columns}
\medskip
\textbf{Questão:} se a leitura do sensor nunca retorna, o controlador consegue atualizar a saída?\par
\medskip
\small Resposta: não nesse ciclo sequencial. A lógica também precisa considerar o tempo de execução.
\end{frame}
